\documentclass[10pt,a4paper]{amsart}
\usepackage{amsmath,amssymb,amsthm,graphicx}
% Retained minimal preamble for the Stacks Project chapter more-morphisms.tex (More on Morphisms).
% Verbatim excerpt of the chapter preamble: Xy-pic setup (native preamble lines 14-18) and the native
% theorem declarations with the mathematical macros (native preamble lines 154-217). Copied mechanically;
% nothing added, removed or redefined. The wrapper class and the amsmath/amssymb/amsthm/graphicx base
% packages are supplied by the assembler, so the class/externaldocument part of the native preamble is omitted.
% For commutative diagrams we use Xy-pic
\usepackage[all]{xy}

% We use 2cell for 2-commutative diagrams.
\xyoption{2cell}
% Theorem environments.
%
\theoremstyle{plain}
\newtheorem{theorem}[subsection]{Theorem}
\newtheorem{proposition}[subsection]{Proposition}
\newtheorem{lemma}[subsection]{Lemma}

\theoremstyle{definition}
\newtheorem{definition}[subsection]{Definition}
\newtheorem{example}[subsection]{Example}
\newtheorem{exercise}[subsection]{Exercise}
\newtheorem{situation}[subsection]{Situation}

\theoremstyle{remark}
\newtheorem{remark}[subsection]{Remark}
\newtheorem{remarks}[subsection]{Remarks}

\numberwithin{equation}{subsection}

% Macros
%
\def\lim{\mathop{\mathrm{lim}}\nolimits}
\def\colim{\mathop{\mathrm{colim}}\nolimits}
\def\Spec{\mathop{\mathrm{Spec}}}
\def\Hom{\mathop{\mathrm{Hom}}\nolimits}
\def\Ext{\mathop{\mathrm{Ext}}\nolimits}
\def\SheafHom{\mathop{\mathcal{H}\!\mathit{om}}\nolimits}
\def\SheafExt{\mathop{\mathcal{E}\!\mathit{xt}}\nolimits}
\def\Sch{\mathit{Sch}}
\def\Mor{\mathop{\mathrm{Mor}}\nolimits}
\def\Ob{\mathop{\mathrm{Ob}}\nolimits}
\def\Sh{\mathop{\mathit{Sh}}\nolimits}
\def\NL{\mathop{N\!L}\nolimits}
\def\CH{\mathop{\mathrm{CH}}\nolimits}
\def\proetale{{pro\text{-}\acute{e}tale}}
\def\etale{{\acute{e}tale}}
\def\QCoh{\mathit{QCoh}}
\def\Ker{\mathop{\mathrm{Ker}}}
\def\Im{\mathop{\mathrm{Im}}}
\def\Coker{\mathop{\mathrm{Coker}}}
\def\Coim{\mathop{\mathrm{Coim}}}

% Boxtimes
%
\DeclareMathSymbol{\boxtimes}{\mathbin}{AMSa}{"02}

%
% Macros for moduli stacks/spaces
%
\def\QCohstack{\mathcal{QC}\!\mathit{oh}}
\def\Cohstack{\mathcal{C}\!\mathit{oh}}
\def\Spacesstack{\mathcal{S}\!\mathit{paces}}
\def\Quotfunctor{\mathrm{Quot}}
\def\Hilbfunctor{\mathrm{Hilb}}
\def\Curvesstack{\mathcal{C}\!\mathit{urves}}
\def\Polarizedstack{\mathcal{P}\!\mathit{olarized}}
\def\Complexesstack{\mathcal{C}\!\mathit{omplexes}}
% \Pic is the operator that assigns to X its picard group, usage \Pic(X)
% \Picardstack_{X/B} denotes the Picard stack of X over B
% \Picardfunctor_{X/B} denotes the Picard functor of X over B
\def\Pic{\mathop{\mathrm{Pic}}\nolimits}
\def\Picardstack{\mathcal{P}\!\mathit{ic}}
\def\Picardfunctor{\mathrm{Pic}}
\def\Deformationcategory{\mathcal{D}\!\mathit{ef}}
\usepackage{cleveref}
\crefname{originaldefinition}{original definition}{original definitions}
\Crefname{originaldefinition}{Original definition}{Original definitions}
\makeatletter
\newlabel{definition-relative-pseudo-coherence}{{09UI}{source}{}{sourcepointer-1}{}}
\newlabel{definition-relative-pseudo-coherence@cref}{{[originaldefinition][0][]09UI}{[1][0][]source}{}{}{}}
\makeatother
\crefname{originallemma}{original lemma}{original lemmas}
\Crefname{originallemma}{Original lemma}{Original lemmas}
\makeatletter
\newlabel{lemma-glue-relative-pseudo-coherent}{{09VE}{source}{}{sourcepointer-2}{}}
\newlabel{lemma-glue-relative-pseudo-coherent@cref}{{[originallemma][0][]09VE}{[1][0][]source}{}{}{}}
\makeatother
\crefname{originallemma}{original lemma}{original lemmas}
\Crefname{originallemma}{Original lemma}{Original lemmas}
\makeatletter
\newlabel{lemma-localize-relative-pseudo-coherent}{{09VD}{source}{}{sourcepointer-3}{}}
\newlabel{lemma-localize-relative-pseudo-coherent@cref}{{[originallemma][0][]09VD}{[1][0][]source}{}{}{}}
\makeatother
\crefname{originallemma}{original lemma}{original lemmas}
\Crefname{originallemma}{Original lemma}{Original lemmas}
\makeatletter
\newlabel{lemma-relatively-pseudo-coherent}{{09VC}{source}{}{sourcepointer-4}{}}
\newlabel{lemma-relatively-pseudo-coherent@cref}{{[originallemma][0][]09VC}{[1][0][]source}{}{}{}}
\makeatother
\crefname{originalreference}{original reference}{original references}
\Crefname{originalreference}{Original reference}{Original references}
\makeatletter
\newlabel{schemes-lemma-standard-open-two-affines}{{01IW}{source}{}{sourcepointer-5}{}}
\newlabel{schemes-lemma-standard-open-two-affines@cref}{{[originalreference][0][]01IW}{[1][0][]source}{}{}{}}
\makeatother
\title{Relative pseudo-coherence is checked on affine pairs}
\author{The Stacks Project Authors}
\date{stacks-project:more-morphisms.tex:local-snapshot-20261011; GNU FDL 1.2 or later}
\begin{document}
\maketitle
\noindent\textit{Editorial note.} Counted span is the author's own proof environment, source lines 16853-16887, immediately following the statement at 16836-16851.\par\medskip
\noindent\textit{Editorial note.} One cross-chapter pointer (Schemes) plus four same-chapter statements are recorded in assembly.external\_references.\par\medskip
\noindent\textit{Editorial note.} The statement's condition on affine pairs is retained verbatim in the statement span.\par\medskip
\noindent\textit{Allowed standard prerequisites (editorial).}\begin{itemize}
\item The affine-pair criterion for one embedding. chapter lemma lemma-relatively-pseudo-coherent, cited in the statement and repeatedly in the proof. Supplies the equivalent conditions that the affine-local formulation refers to.
\item Localisation of relative pseudo-coherence. chapter lemma lemma-localize-relative-pseudo-coherent, cited at proof lines 16880-16883. Moves the condition along the standard affine open U' inside the smaller pair.
\item Glueing relative pseudo-coherence. chapter lemma lemma-glue-relative-pseudo-coherent, cited in the proof. Reassembles the local verifications into the global statement.
\item Standard affine opens inside two affine opens. Schemes, lemma schemes-lemma-standard-open-two-affines, cited in the proof. Produces the standard affine open contained in the intersections appearing in the localisation argument.
\end{itemize}\medskip
\ifdefined\hypertarget\hypertarget{sourcepointer-1}{}\fi
\noindent\textit{Source reference.} Original definition 09UI: sources/stacks-more-morphisms/pinned-original-source.tex, source lines 16675--16675. Navigation pointer only: the counted proof cites this statement by \textbackslash{}ref as a standing input; the deduction the counted proof has to make is carried out inside the counted span, and no part of that deduction is delegated to this pointer.\par\medskip
\ifdefined\hypertarget\hypertarget{sourcepointer-2}{}\fi
\noindent\textit{Source reference.} Original lemma 09VE: sources/stacks-more-morphisms/pinned-original-source.tex, source lines 16789--16789. Navigation pointer only: the counted proof cites this statement by \textbackslash{}ref as a standing input; the deduction the counted proof has to make is carried out inside the counted span, and no part of that deduction is delegated to this pointer.\par\medskip
\ifdefined\hypertarget\hypertarget{sourcepointer-3}{}\fi
\noindent\textit{Source reference.} Original lemma 09VD: sources/stacks-more-morphisms/pinned-original-source.tex, source lines 16738--16738. Navigation pointer only: the counted proof cites this statement by \textbackslash{}ref as a standing input; the deduction the counted proof has to make is carried out inside the counted span, and no part of that deduction is delegated to this pointer.\par\medskip
\ifdefined\hypertarget\hypertarget{sourcepointer-4}{}\fi
\noindent\textit{Source reference.} Original lemma 09VC: sources/stacks-more-morphisms/pinned-original-source.tex, source lines 16610--16610. Navigation pointer only: the counted proof cites this statement by \textbackslash{}ref as a standing input; the deduction the counted proof has to make is carried out inside the counted span, and no part of that deduction is delegated to this pointer.\par\medskip
\ifdefined\hypertarget\hypertarget{sourcepointer-5}{}\fi
\noindent\textit{Source reference.} Original reference 01IW: focused-seeds/geometry/S9-morph-rel-pseudo-coherence-characterize/ref-tags.txt, source lines 5--5. Navigation pointer only: the counted proof cites this statement by \textbackslash{}ref as a standing input; the deduction the counted proof has to make is carried out inside the counted span, and no part of that deduction is delegated to this pointer.\par\medskip
\begin{lemma}
\label{lemma-relative-pseudo-coherence-characterize}
Let $f : X \to S$ be a morphism of schemes which is locally of finite type.
Let $E$ be an object of $D(\mathcal{O}_X)$.
Fix $m \in \mathbf{Z}$. The following are equivalent
\begin{enumerate}
\item $E$ is $m$-pseudo-coherent relative to $S$,
\item for every affine opens $U \subset X$ and $V \subset S$
with $f(U) \subset V$ the equivalent conditions of
Lemma \ref{lemma-relatively-pseudo-coherent}
are satisfied for the pair $(U \to V, E|_U)$.
\end{enumerate}
Moreover, if this is true, then for every open subschemes
$U \subset X$ and $V \subset S$ with $f(U) \subset V$
the restriction $E|_U$ is $m$-pseudo-coherent relative to $V$.
\end{lemma}

\begin{proof}
The final statement is clear from the equivalence of (1) and (2).
It is also clear that (2) implies (1). Assume (1) holds.
Let $S = \bigcup V_i$ and $f^{-1}(V_i) = \bigcup U_{ij}$ be
affine open coverings as in
Definition \ref{definition-relative-pseudo-coherence}.
Let $U \subset X$ and $V \subset S$ be as in (2).
By Lemma \ref{lemma-glue-relative-pseudo-coherent}
it suffices to find a standard open covering $U = \bigcup U_k$ of $U$
such that the equivalent conditions of
Lemma \ref{lemma-relatively-pseudo-coherent}
are satisfied for the pairs $(U_k \to V, E|_{U_k})$.
In other words, for every $u \in U$ it suffices
to find a standard affine open $u \in U' \subset U$ such that
the equivalent conditions of
Lemma \ref{lemma-relatively-pseudo-coherent}
are satisfied for the pair $(U' \to V, E|_{U'})$.
Pick $i$ such that $f(u) \in V_i$ and then pick $j$ such that
$u \in U_{ij}$. By
Schemes, Lemma \ref{schemes-lemma-standard-open-two-affines}
we can find $v \in V' \subset V \cap V_i$ which is standard affine
open in $V'$ and $V_i$. Then $f^{-1}V'  \cap U$, resp.\ $f^{-1}V' \cap U_{ij}$
are standard affine opens of $U$, resp.\ $U_{ij}$.
Applying the lemma again we can find
$u \in U' \subset f^{-1}V' \cap U \cap U_{ij}$ which is standard affine
open in both $f^{-1}V'  \cap U$ and $f^{-1}V' \cap U_{ij}$.
Thus $U'$ is also a standard affine open of $U$ and $U_{ij}$.
By Lemma \ref{lemma-localize-relative-pseudo-coherent}
the assumption that the equivalent conditions of
Lemma \ref{lemma-relatively-pseudo-coherent}
are satisfied for the pair $(U_{ij} \to V_i, E|_{U_{ij}})$
implies that the equivalent conditions of
Lemma \ref{lemma-relatively-pseudo-coherent}
are satisfied for the pair $(U' \to V, E|_{U'})$.
\end{proof}

\end{document}
